\chapter{Conclusão} \label{cpt:conclus}
Este trabalho apresentou uma proposta de abordagem para contribuir com o desenvolvimento e pesquisa na área da saúde. Inicialmente, desenvolvemos um conjunto de dados contendo informações reais sobre a quantidade de leitos hospitalares, classificados por especialidade e tipo, nos hospitais de Porto Alegre. Além disso, incluímos nesses dados informações geográficas de latitude e longitude de cada hospital, tornando o conjunto útil para uma variedade de aplicações.

Em complemento, propusemos o LODUS-Health, uma ferramenta que permite simular a ocupação de leitos hospitalares na cidade de Porto Alegre oferecendo uma interface gráfica intuitiva, facilitando a interação dos usuários e possibilitando parametrizações customizadas. Além disso, ela integra os dados do conjunto previamente criado e permite a geração de dados médicos específicos para cada especialidade de leito selecionada a partir dos resultados gerados na simulação.

Os resultados obtidos indicam que o LODUS-Health pode ser uma ferramenta valiosa para os gestores hospitalares e urbanos da cidade, no auxilio na alocação eficiente de recursos e na tomada de decisões estratégicas, como a inclusão de novos hospitais em determinadas áreas.

Apesar das contribuições significativas, é importante reconhecer as limitações deste estudo. O modelo de simulação apresentado requer um certo nível de conhecimento prévio sobre o fluxo de ocupação de leitos hospitalares, o que pode limitar sua aplicabilidade em alguns contextos. Além disso, a necessidade de atualização dos dados sobre os leitos hospitalares representa um desafio, apesar de acreditarmos que mudanças significativas na quantidade de leitos não sejam realizadas com frequência.